% Options for packages loaded elsewhere
\PassOptionsToPackage{unicode}{hyperref}
\PassOptionsToPackage{hyphens}{url}
%
\documentclass[
  ignorenonframetext,
]{beamer}
\usepackage{pgfpages}
\setbeamertemplate{caption}[numbered]
\setbeamertemplate{caption label separator}{: }
\setbeamercolor{caption name}{fg=normal text.fg}
\beamertemplatenavigationsymbolsempty
% Prevent slide breaks in the middle of a paragraph
\widowpenalties 1 10000
\raggedbottom
\setbeamertemplate{part page}{
  \centering
  \begin{beamercolorbox}[sep=16pt,center]{part title}
    \usebeamerfont{part title}\insertpart\par
  \end{beamercolorbox}
}
\setbeamertemplate{section page}{
  \centering
  \begin{beamercolorbox}[sep=12pt,center]{part title}
    \usebeamerfont{section title}\insertsection\par
  \end{beamercolorbox}
}
\setbeamertemplate{subsection page}{
  \centering
  \begin{beamercolorbox}[sep=8pt,center]{part title}
    \usebeamerfont{subsection title}\insertsubsection\par
  \end{beamercolorbox}
}
\AtBeginPart{
  \frame{\partpage}
}
\AtBeginSection{
  \ifbibliography
  \else
    \frame{\sectionpage}
  \fi
}
\AtBeginSubsection{
  \frame{\subsectionpage}
}
\usepackage{amsmath,amssymb}
\usepackage{iftex}
\ifPDFTeX
  \usepackage[T1]{fontenc}
  \usepackage[utf8]{inputenc}
  \usepackage{textcomp} % provide euro and other symbols
\else % if luatex or xetex
  \usepackage{unicode-math} % this also loads fontspec
  \defaultfontfeatures{Scale=MatchLowercase}
  \defaultfontfeatures[\rmfamily]{Ligatures=TeX,Scale=1}
\fi
\usepackage{lmodern}
\ifPDFTeX\else
  % xetex/luatex font selection
\fi
% Use upquote if available, for straight quotes in verbatim environments
\IfFileExists{upquote.sty}{\usepackage{upquote}}{}
\IfFileExists{microtype.sty}{% use microtype if available
  \usepackage[]{microtype}
  \UseMicrotypeSet[protrusion]{basicmath} % disable protrusion for tt fonts
}{}
\makeatletter
\@ifundefined{KOMAClassName}{% if non-KOMA class
  \IfFileExists{parskip.sty}{%
    \usepackage{parskip}
  }{% else
    \setlength{\parindent}{0pt}
    \setlength{\parskip}{6pt plus 2pt minus 1pt}}
}{% if KOMA class
  \KOMAoptions{parskip=half}}
\makeatother
\usepackage{xcolor}
\newif\ifbibliography
\usepackage{color}
\usepackage{fancyvrb}
\newcommand{\VerbBar}{|}
\newcommand{\VERB}{\Verb[commandchars=\\\{\}]}
\DefineVerbatimEnvironment{Highlighting}{Verbatim}{commandchars=\\\{\}}
% Add ',fontsize=\small' for more characters per line
\usepackage{framed}
\definecolor{shadecolor}{RGB}{248,248,248}
\newenvironment{Shaded}{\begin{snugshade}}{\end{snugshade}}
\newcommand{\AlertTok}[1]{\textcolor[rgb]{0.94,0.16,0.16}{#1}}
\newcommand{\AnnotationTok}[1]{\textcolor[rgb]{0.56,0.35,0.01}{\textbf{\textit{#1}}}}
\newcommand{\AttributeTok}[1]{\textcolor[rgb]{0.13,0.29,0.53}{#1}}
\newcommand{\BaseNTok}[1]{\textcolor[rgb]{0.00,0.00,0.81}{#1}}
\newcommand{\BuiltInTok}[1]{#1}
\newcommand{\CharTok}[1]{\textcolor[rgb]{0.31,0.60,0.02}{#1}}
\newcommand{\CommentTok}[1]{\textcolor[rgb]{0.56,0.35,0.01}{\textit{#1}}}
\newcommand{\CommentVarTok}[1]{\textcolor[rgb]{0.56,0.35,0.01}{\textbf{\textit{#1}}}}
\newcommand{\ConstantTok}[1]{\textcolor[rgb]{0.56,0.35,0.01}{#1}}
\newcommand{\ControlFlowTok}[1]{\textcolor[rgb]{0.13,0.29,0.53}{\textbf{#1}}}
\newcommand{\DataTypeTok}[1]{\textcolor[rgb]{0.13,0.29,0.53}{#1}}
\newcommand{\DecValTok}[1]{\textcolor[rgb]{0.00,0.00,0.81}{#1}}
\newcommand{\DocumentationTok}[1]{\textcolor[rgb]{0.56,0.35,0.01}{\textbf{\textit{#1}}}}
\newcommand{\ErrorTok}[1]{\textcolor[rgb]{0.64,0.00,0.00}{\textbf{#1}}}
\newcommand{\ExtensionTok}[1]{#1}
\newcommand{\FloatTok}[1]{\textcolor[rgb]{0.00,0.00,0.81}{#1}}
\newcommand{\FunctionTok}[1]{\textcolor[rgb]{0.13,0.29,0.53}{\textbf{#1}}}
\newcommand{\ImportTok}[1]{#1}
\newcommand{\InformationTok}[1]{\textcolor[rgb]{0.56,0.35,0.01}{\textbf{\textit{#1}}}}
\newcommand{\KeywordTok}[1]{\textcolor[rgb]{0.13,0.29,0.53}{\textbf{#1}}}
\newcommand{\NormalTok}[1]{#1}
\newcommand{\OperatorTok}[1]{\textcolor[rgb]{0.81,0.36,0.00}{\textbf{#1}}}
\newcommand{\OtherTok}[1]{\textcolor[rgb]{0.56,0.35,0.01}{#1}}
\newcommand{\PreprocessorTok}[1]{\textcolor[rgb]{0.56,0.35,0.01}{\textit{#1}}}
\newcommand{\RegionMarkerTok}[1]{#1}
\newcommand{\SpecialCharTok}[1]{\textcolor[rgb]{0.81,0.36,0.00}{\textbf{#1}}}
\newcommand{\SpecialStringTok}[1]{\textcolor[rgb]{0.31,0.60,0.02}{#1}}
\newcommand{\StringTok}[1]{\textcolor[rgb]{0.31,0.60,0.02}{#1}}
\newcommand{\VariableTok}[1]{\textcolor[rgb]{0.00,0.00,0.00}{#1}}
\newcommand{\VerbatimStringTok}[1]{\textcolor[rgb]{0.31,0.60,0.02}{#1}}
\newcommand{\WarningTok}[1]{\textcolor[rgb]{0.56,0.35,0.01}{\textbf{\textit{#1}}}}
\usepackage{longtable,booktabs,array}
\usepackage{calc} % for calculating minipage widths
\usepackage{caption}
% Make caption package work with longtable
\makeatletter
\def\fnum@table{\tablename~\thetable}
\makeatother
\setlength{\emergencystretch}{3em} % prevent overfull lines
\providecommand{\tightlist}{%
  \setlength{\itemsep}{0pt}\setlength{\parskip}{0pt}}
\setcounter{secnumdepth}{-\maxdimen} % remove section numbering
\ifLuaTeX
  \usepackage{selnolig}  % disable illegal ligatures
\fi
\usepackage{bookmark}
\IfFileExists{xurl.sty}{\usepackage{xurl}}{} % add URL line breaks if available
\urlstyle{same}
\hypersetup{
  pdftitle={IPWMethod Package Overview},
  pdfauthor={Jiakun},
  hidelinks,
  pdfcreator={LaTeX via pandoc}}

\title{IPWMethod Package Overview}
\author{Jiakun}
\date{2025-12-02}

\begin{document}
\frame{\titlepage}

\begin{frame}{Presentation Overview}
\phantomsection\label{presentation-overview}
\begin{itemize}
\item
  Example datasets (NHANES and NHIS)
\item
  Package

  \begin{itemize}
  \tightlist
  \item
    Function interface
  \item
    Examples
  \item
    Structure of output objects
  \end{itemize}
\item
  Questions

  \begin{itemize}
  \tightlist
  \item
    Pre-calibration
  \item
    Missing data
  \item
    Function structure
  \end{itemize}
\end{itemize}
\end{frame}

\begin{frame}{NHANES (2000-2010 five waves)}
\phantomsection\label{nhanes-2000-2010-five-waves}
\begin{longtable}[]{@{}rlr@{}}
\toprule\noalign{}
Wave & Year & Number of subjects \\
\midrule\noalign{}
\endhead
1 & 2001-2002 & 676 \\
2 & 2003-2004 & 677 \\
3 & 2005-2006 & 632 \\
4 & 2007-2008 & 863 \\
5 & 2009-2010 & 898 \\
\bottomrule\noalign{}
\end{longtable}
\end{frame}

\begin{frame}[fragile]{sc}
\phantomsection\label{sc}
sc is from 2000-2004 NHANES wave\\
sc\_full with missing value (1353 rows)\\
sc\_cc (complete case) without missing value (1082 rows)\\
height and weight are self-reported to mimic PLCO\\
true weights are kept for comparison

\begin{verbatim}
## 'data.frame':    1353 obs. of  12 variables:
##  $ agecat     : Factor w/ 4 levels "1","2","3","4": 2 1 4 2 3 2 4 4 2 4 ...
##  $ marital    : Factor w/ 4 levels "1","2","3","4": 1 3 1 1 2 1 1 1 1 1 ...
##  $ race       : Factor w/ 4 levels "1","2","3","4": 4 3 1 1 2 2 3 1 3 3 ...
##  $ educat     : Factor w/ 5 levels "1","2","3","4",..: 4 1 3 4 5 1 2 5 3 1 ...
##  $ empstat    : Factor w/ 2 levels "0","1": 2 2 2 1 1 2 1 1 2 1 ...
##  $ smoking    : Factor w/ 3 levels "1","2","3": 2 3 2 2 3 2 1 2 3 3 ...
##  $ comorbidity: Factor w/ 2 levels "0","1": 2 2 1 1 2 2 2 2 1 1 ...
##  $ psa        : num  0.4 3.6 1.3 1 1 0.9 2.7 0.5 0.8 1.2 ...
##  $ height     : num  1.6 1.73 1.68 1.85 1.8 ...
##  $ weight     : num  57.6 85.3 80.7 97.5 82.1 ...
##  $ BMI        : num  22.5 28.6 28.7 28.4 25.2 ...
##  $ true_wts   : num  13100 1888 5239 9589 1400 ...
\end{verbatim}
\end{frame}

\begin{frame}[fragile]{sp1}
\phantomsection\label{sp1}
sp1 is from 2005-2010 NHANES wave\\
sp1\_full with missing values (1495 rows)\\
sp1\_cc (complete case) without missing values (1229 rows)\\
self-reported height and weight are not ``measured'' in NHANES

\begin{verbatim}
## 'data.frame':    1495 obs. of  9 variables:
##  $ agecat     : Factor w/ 4 levels "1","2","3","4": 4 4 2 4 4 4 4 2 2 1 ...
##  $ marital    : Factor w/ 4 levels "1","2","3","4": 1 1 1 1 3 1 1 1 1 1 ...
##  $ race       : Factor w/ 4 levels "1","2","3","4": 1 1 2 1 1 1 1 3 1 1 ...
##  $ educat     : Factor w/ 5 levels "1","2","3","4",..: 5 3 3 4 3 4 4 2 5 5 ...
##  $ empstat    : Factor w/ 2 levels "0","1": 2 1 2 1 1 2 2 1 1 2 ...
##  $ smoking    : Factor w/ 3 levels "1","2","3": 1 1 3 3 2 1 2 1 1 1 ...
##  $ comorbidity: Factor w/ 2 levels "0","1": 2 2 2 2 2 2 2 2 2 1 ...
##  $ psa        : num  2.25 2.15 0.91 1.83 6.51 3.63 0.98 2.76 8.68 6.29 ...
##  $ wts_sp1    : num  6592 8744 1853 8141 8488 ...
\end{verbatim}
\end{frame}

\begin{frame}[fragile]{sp2}
\phantomsection\label{sp2}
sp2 is from 1997-2008 NHIS (male 50-75 yrs)\\
sp2\_full with missing values (35610 rows)\\
sp2\_cc (complete case) without missing values (35293 rows) height and
weight are self-reported

\begin{verbatim}
## 'data.frame':    35610 obs. of  11 variables:
##  $ agecat     : Factor w/ 4 levels "1","2","3","4": 3 1 3 1 4 4 1 1 1 1 ...
##  $ marital    : Factor w/ 4 levels "1","2","3","4": 2 1 4 1 1 2 2 3 1 1 ...
##  $ race       : Factor w/ 4 levels "1","2","3","4": 3 3 3 3 3 3 2 2 2 3 ...
##  $ empstat    : Factor w/ 2 levels "0","1": 2 2 1 2 2 1 1 1 1 1 ...
##  $ bmicat     : Factor w/ 5 levels "0","1","2","3",..: 2 2 3 2 4 3 2 3 5 4 ...
##  $ smoking    : Factor w/ 3 levels "1","2","3": 1 1 2 3 2 2 3 3 2 1 ...
##  $ comorbidity: Factor w/ 2 levels "0","1": 2 2 2 1 1 1 2 2 2 2 ...
##  $ height     : num  1.63 1.85 1.73 1.68 1.63 ...
##  $ weight     : num  64 81.6 74.8 68 79.4 ...
##  $ BMI        : num  24.2 23.7 25.1 24.2 30 ...
##  $ wts_sp2    : int  2466 8271 2913 9368 3401 1701 2253 2255 5829 1452 ...
\end{verbatim}
\end{frame}

\begin{frame}[fragile]{IPWMethod Package}
\phantomsection\label{ipwmethod-package}
\begin{itemize}
\tightlist
\item
  One main function IPWM\\
\item
  Support three single-reference methods - ALP, CLW, and calibration\\
\item
  One multi-reference method that performs calibration with optional
  pre-calibration
\end{itemize}

\begin{Shaded}
\begin{Highlighting}[]
\NormalTok{IPWM }\OtherTok{\textless{}{-}} \ControlFlowTok{function}\NormalTok{(}
\NormalTok{    sc,                          }\CommentTok{\# sc: convenience sample data.frame}
\NormalTok{    sp,                          }\CommentTok{\# sp: data.frame (One) or list of data.frames (Multi)}
\NormalTok{    y,                           }\CommentTok{\# y: outcome variable in sc}
\NormalTok{    weight,                      }\CommentTok{\# weight: string (One) or vector (Multi)}
    \AttributeTok{p\_formula =} \ConstantTok{NULL}\NormalTok{,            }\CommentTok{\# p\_formula: formula (One) or list of formulas (Multi)}
    \AttributeTok{method =} \ConstantTok{NULL}\NormalTok{,               }\CommentTok{\# method: "ALP", "CLW", "calibration", or "multi"}
    \AttributeTok{zcol =} \ConstantTok{NULL}\NormalTok{,                 }\CommentTok{\# zcol: domain variable}
    \AttributeTok{cali =} \ConstantTok{TRUE}\NormalTok{,                 }\CommentTok{\# cali: whether to pre{-}calibrate}
    
    
    \AttributeTok{maxit =} \DecValTok{20}\NormalTok{,                  }\CommentTok{\# maxit: maximum Newton–Raphson iterations}
    \AttributeTok{tol =} \FloatTok{1e{-}4}\NormalTok{,                  }\CommentTok{\# tol: convergence tolerance}
    \AttributeTok{verbose =} \ConstantTok{FALSE}              \CommentTok{\# verbose: print messages}
\NormalTok{)}
\end{Highlighting}
\end{Shaded}
\end{frame}

\begin{frame}[fragile]{One reference method}
\phantomsection\label{one-reference-method}
standard use of function

\begin{Shaded}
\begin{Highlighting}[]
\FunctionTok{IPWM}\NormalTok{(}\AttributeTok{sc=}\NormalTok{sc\_cc, }\AttributeTok{sp=}\NormalTok{sp1\_cc, }\AttributeTok{y=}\StringTok{\textquotesingle{}psa\textquotesingle{}}\NormalTok{, }\AttributeTok{weight =} \StringTok{\textquotesingle{}wts\_sp1\textquotesingle{}}\NormalTok{, }\AttributeTok{method =} \StringTok{"ALP"}\NormalTok{, }
     \AttributeTok{p\_formula =} \SpecialCharTok{\textasciitilde{}}\NormalTok{ agecat }\SpecialCharTok{+}\NormalTok{ marital }\SpecialCharTok{+}\NormalTok{ race }\SpecialCharTok{+}\NormalTok{ educat }\SpecialCharTok{+}\NormalTok{ psa )}
\end{Highlighting}
\end{Shaded}

\begin{verbatim}
## Call:
## IPWM(sc = sc_cc, sp = sp1_cc, y = "psa", weight = "wts_sp1", 
##     p_formula = ~agecat + marital + race + educat + psa, method = "ALP")
## 
## Pseudo-weighted Estimators:
##   Mean:         1.943621 
##   Std. Error:   0.112599 
##   95% CI:      [1.722928, 2.164315]
\end{verbatim}
\end{frame}

\begin{frame}[fragile]{One reference method}
\phantomsection\label{one-reference-method-1}
output summary

\begin{Shaded}
\begin{Highlighting}[]
\NormalTok{result }\OtherTok{=} \FunctionTok{IPWM}\NormalTok{(}\AttributeTok{sc=}\NormalTok{sc\_cc, }\AttributeTok{sp=}\NormalTok{sp1\_cc, }\AttributeTok{y=}\StringTok{\textquotesingle{}psa\textquotesingle{}}\NormalTok{, }\AttributeTok{weight =} \StringTok{\textquotesingle{}wts\_sp1\textquotesingle{}}\NormalTok{, }\AttributeTok{method =} \StringTok{"ALP"}\NormalTok{, }
              \AttributeTok{p\_formula =} \SpecialCharTok{\textasciitilde{}}\NormalTok{ agecat }\SpecialCharTok{+}\NormalTok{ marital }\SpecialCharTok{+}\NormalTok{ race }\SpecialCharTok{+}\NormalTok{ educat }\SpecialCharTok{+}\NormalTok{ psa )}
\FunctionTok{summary}\NormalTok{(result)}
\end{Highlighting}
\end{Shaded}

\begin{verbatim}
## Call:
## IPWM(sc = sc_cc, sp = sp1_cc, y = "psa", weight = "wts_sp1", 
##     p_formula = ~agecat + marital + race + educat + psa, method = "ALP")
## 
## Method: One reference ALP 
## 
## Participation model involves the following variables:
## agecat2 agecat3 agecat4 marital2 marital3 marital4 race2 race3 race4 educat2 educat3 educat4 educat5 psa 
## 
## Unweighted Estimators summary:
##   Mean:                    1.969316
##   Std. Error:              0.006298
##   95% CI:                [1.813772,2.124860]
## 
## Pseudo-weighted (ALP) Estimators summary:
##   Mean:                    1.943621
##   Std. Error:              0.012679
##   95% CI:                [1.722928,2.164315]
##   (Newton–Raphson iterations: 7)
## 
## Selection model coefficients:
##  (Intercept)    agecat2    agecat3    agecat4   marital2   marital3   marital4
##      -9.2516     0.5173     0.6818     1.1167    -0.2769    -0.1737    -0.3280
##       race2      race3      race4    educat2    educat3    educat4    educat5
##      1.1085     1.5185    -0.0669    -0.3535    -0.5621    -0.3827    -0.5674
##         psa
##     -0.0288
\end{verbatim}
\end{frame}

\begin{frame}[fragile]{Multi-reference method}
\phantomsection\label{multi-reference-method}
\begin{Shaded}
\begin{Highlighting}[]
\NormalTok{result }\OtherTok{=} \FunctionTok{IPWM}\NormalTok{(}\AttributeTok{sc=}\NormalTok{sc\_cc, }\AttributeTok{sp=}\FunctionTok{list}\NormalTok{(sp1\_cc,sp2\_cc), }\AttributeTok{y=}\StringTok{\textquotesingle{}psa\textquotesingle{}}\NormalTok{, }
              \AttributeTok{weight =} \FunctionTok{c}\NormalTok{(}\StringTok{\textquotesingle{}wts\_sp1\textquotesingle{}}\NormalTok{,}\StringTok{\textquotesingle{}wts\_sp2\textquotesingle{}}\NormalTok{),  }\AttributeTok{cali =} \ConstantTok{TRUE}\NormalTok{,}
              \AttributeTok{p\_formula =} \FunctionTok{list}\NormalTok{( }\SpecialCharTok{\textasciitilde{}}\NormalTok{ psa ,}
                                \SpecialCharTok{\textasciitilde{}}\NormalTok{ agecat }\SpecialCharTok{+}\NormalTok{ marital }\SpecialCharTok{+}\NormalTok{ race }\SpecialCharTok{+}\NormalTok{ empstat }\SpecialCharTok{+}\NormalTok{ smoking }\SpecialCharTok{+}\NormalTok{ comorbidity }\SpecialCharTok{+}\NormalTok{  BMI))}
\FunctionTok{summary}\NormalTok{(result)}
\end{Highlighting}
\end{Shaded}

\begin{verbatim}
## Call:
## IPWM(sc = sc_cc, sp = list(sp1_cc, sp2_cc), y = "psa", weight = c("wts_sp1", 
##     "wts_sp2"), p_formula = list(~psa, ~agecat + marital + race + 
##     empstat + smoking + comorbidity + BMI), cali = TRUE)
## 
## Method: Multi reference calibration
## 
## Pre-calibration summary:
## Pre-calibration reference: sp[[2]] (n = 35293)
## Pre-calibration for sp[[1]] done on: agecat2, agecat3, agecat4, comorbidity1, empstat1, marital2, marital3, marital4, race2, race3, race4, smoking2, smoking3 and survey weights total
## 
## Reference samples summary:
## Order of samples by size (largest to smallest):
##   sp[[2]] (n = 35293)
##   sp[[1]] (n = 1240)
## Shared variables in sp[[2]]:
##   agecat2, agecat3, agecat4, marital2, marital3, marital4, race2, race3, race4, empstat1, smoking2, smoking3, comorbidity1, BMI
## Shared variables in sp[[1]]:
##   psa
## Variables used for calculation in sp[[2]]:
##   agecat2, agecat3, agecat4, marital2, marital3, marital4, race2, race3, race4, empstat1, smoking2, smoking3, comorbidity1, BMI
## Variables used for calculation in sp[[1]]:
##   psa
## 
## Participation model involves the following variables:
## agecat2 agecat3 agecat4 marital2 marital3 marital4 race2 race3 race4 empstat1 smoking2 smoking3 comorbidity1 BMI psa 
## 
## Unweighted Estimators summary:
##   Mean:                    1.969316
##   Std. Error:              0.006298
##   95% CI:                [1.813772,2.124860]
## 
## Pseudo-weighted (multi) Estimators summary:
##   Mean:                    2.018745
##   Std. Error:              0.014022
##   95% CI:                [1.786653,2.250836]
##   (Newton–Raphson iterations: 6)
## 
## Selection model coefficients:
##  (Intercept)    agecat2    agecat3    agecat4   marital2   marital3   marital4
##     -10.7500     0.4948     0.6007     0.7787    -0.2018    -0.0428     0.0988
##       race2      race3      race4   empstat1   smoking2   smoking3 comorbidity1
##      1.1648     1.5417    -0.1901    -0.1510     0.2096     0.1866       0.1452
##         BMI        psa
##     -0.0859    -0.0297
\end{verbatim}
\end{frame}

\begin{frame}[fragile]{Multi-reference method}
\phantomsection\label{multi-reference-method-1}
\begin{Shaded}
\begin{Highlighting}[]
\NormalTok{result }\OtherTok{=} \FunctionTok{IPWM}\NormalTok{(}\AttributeTok{sc=}\NormalTok{sc\_full, }\AttributeTok{sp=}\FunctionTok{list}\NormalTok{(sp1\_full,sp2\_full), }\AttributeTok{y=}\StringTok{\textquotesingle{}psa\textquotesingle{}}\NormalTok{, }
              \AttributeTok{weight =} \FunctionTok{c}\NormalTok{(}\StringTok{\textquotesingle{}wts\_sp1\textquotesingle{}}\NormalTok{,}\StringTok{\textquotesingle{}wts\_sp2\textquotesingle{}}\NormalTok{), }\AttributeTok{cali =} \ConstantTok{TRUE}\NormalTok{,}
              \AttributeTok{p\_formula =} \FunctionTok{list}\NormalTok{( }\SpecialCharTok{\textasciitilde{}}\NormalTok{ psa ,}
                                \SpecialCharTok{\textasciitilde{}}\NormalTok{ agecat }\SpecialCharTok{+}\NormalTok{ marital }\SpecialCharTok{+}\NormalTok{ race }\SpecialCharTok{+}\NormalTok{ empstat }\SpecialCharTok{+}\NormalTok{ smoking }\SpecialCharTok{+}\NormalTok{ comorbidity }\SpecialCharTok{+}\NormalTok{  BMI))}
\FunctionTok{summary}\NormalTok{(result)}
\end{Highlighting}
\end{Shaded}

\begin{verbatim}
## Call:
## IPWM(sc = sc_full, sp = list(sp1_full, sp2_full), y = "psa", 
##     weight = c("wts_sp1", "wts_sp2"), p_formula = list(~psa, 
##         ~agecat + marital + race + empstat + smoking + comorbidity + 
##             BMI), cali = TRUE)
## 
## Method: Multi reference calibration
## 
## Missing value summary:
## 272 NA detected in p_formula variables in sc:  - psa  - BMI
## Removed 271 rows from sc due to NA in p_formula variables.
## 255 NA detected in p_formula variables in sp[[1]]:  - psa
## Removed 255 rows from sp[[1]] due to NA in p_formula variables.
## 317 NA detected in p_formula variables in sp[[2]]:  - BMI
## Removed 317 rows from sp[[2]] due to NA in p_formula variables.
## 
## Pre-calibration summary:
## Pre-calibration reference: sp[[2]] (n = 35293)
## Pre-calibration for sp[[1]] done on: agecat2, agecat3, agecat4, comorbidity1, empstat1, marital2, marital3, marital4, race2, race3, race4, smoking2, smoking3 and survey weights total
## 
## Reference samples summary:
## Order of samples by size (largest to smallest):
##   sp[[2]] (n = 35293)
##   sp[[1]] (n = 1240)
## Shared variables in sp[[2]]:
##   agecat2, agecat3, agecat4, marital2, marital3, marital4, race2, race3, race4, empstat1, smoking2, smoking3, comorbidity1, BMI
## Shared variables in sp[[1]]:
##   psa
## Variables used for calculation in sp[[2]]:
##   agecat2, agecat3, agecat4, marital2, marital3, marital4, race2, race3, race4, empstat1, smoking2, smoking3, comorbidity1, BMI
## Variables used for calculation in sp[[1]]:
##   psa
## 
## Participation model involves the following variables:
## agecat2 agecat3 agecat4 marital2 marital3 marital4 race2 race3 race4 empstat1 smoking2 smoking3 comorbidity1 BMI psa 
## 
## Unweighted Estimators summary:
##   Mean:                    1.969316
##   Std. Error:              0.006298
##   95% CI:                [1.813772,2.124860]
## 
## Pseudo-weighted (multi) Estimators summary:
##   Mean:                    2.018745
##   Std. Error:              0.014022
##   95% CI:                [1.786653,2.250836]
##   (Newton–Raphson iterations: 6)
## 
## Selection model coefficients:
##  (Intercept)    agecat2    agecat3    agecat4   marital2   marital3   marital4
##     -10.7500     0.4948     0.6007     0.7787    -0.2018    -0.0428     0.0988
##       race2      race3      race4   empstat1   smoking2   smoking3 comorbidity1
##      1.1648     1.5417    -0.1901    -0.1510     0.2096     0.1866       0.1452
##         BMI        psa
##     -0.0859    -0.0297
\end{verbatim}
\end{frame}

\begin{frame}{Question: Pre-calibration in three samples}
\phantomsection\label{question-pre-calibration-in-three-samples}
\end{frame}

\begin{frame}{Question: full vs complete case data}
\phantomsection\label{question-full-vs-complete-case-data}
\begin{itemize}
\tightlist
\item
  complete case data justification:

  \begin{itemize}
  \tightlist
  \item
    follow the paper\\
  \item
    multi imputation is out of scope for now
  \end{itemize}
\item
  complete case data problem:

  \begin{itemize}
  \tightlist
  \item
    some people don't get pseudo weights assigned when outcome is
    measured
  \item
    consider basic imputation of pseudo weights?
  \end{itemize}
\item
  mimic how lm handles NA?

  \begin{itemize}
  \tightlist
  \item
    lm(y \textasciitilde{} x, data = df, na.action = na.omit) \# Default
    (remove rows)\\
  \item
    lm(y \textasciitilde{} x, data = df, na.action = na.fail) \# Stops
    with error if any NA exists\\
  \item
    fit \textless- lm(y \textasciitilde{} x1 + x2, data = df)\\
  \item
    fit\$na.action \# shows which rows are dropped, not printed in
    summary
  \end{itemize}
\end{itemize}
\end{frame}

\begin{frame}{Question - Separation of functions}
\phantomsection\label{question---separation-of-functions}
\begin{itemize}
\item
  now one big function that returns mean outcome
\item
  separation of functions to give pseudo-weights based on methods and
  mean outcome separately (fixed vs flexiable)

  \begin{itemize}
  \tightlist
  \item
    one function returns a dataset with the pseudo-weights
  \item
    another function returns the outcome
  \item
    make future work easier (regression/hypothesis testing)
  \end{itemize}
\end{itemize}
\end{frame}

\end{document}
